\chapter{Convergence and area decay for periodic graphs}\label{sec:periodic-graph-convergence}
The curvature estimate from \cite{EH89}, Proposition 4.4, gives a useful compact-domain consequence. Throughout this section, the flow is assumed smooth for every $t\geq0$ on a fixed flat torus $\mathbb T^n$. We prove its convergence, including quantitative area decay. These conclusions do not replace a proof of global existence.

\begin{lemma}[Converting graph curvature to Euclidean Hessian bounds]\label{lem:graph-curvature-hessian-comparison}
For a smooth graph with $v=\sqrt{1+|Du|^2}$, its metric and scalar second fundamental form satisfy
\[
 g_{ij}=\delta_{ij}+u_i u_j,\qquad h_{ij}=-v^{-1}u_{ij},\qquad
 v^{-6}|D^2u|^2\leq|A|^2\leq v^{-2}|D^2u|^2.
\]
The Hessian norm on the outer terms is Euclidean. In particular $|D^2u|\leq v^3|A|$.
\end{lemma}
\begin{proof}
The metric follows by differentiating the graph map. Differentiating $\langle\nu,F_j\rangle=0$ gives $h_{ij}=-\langle\nu,F_{ij}\rangle=-u_{ij}/v$. In a Euclidean orthonormal basis diagonalizing $g^{-1}$, its eigenvalues lie between $v^{-2}$ and one. The squared tensor norm is $v^{-2}\sum_{i,j}\gamma_i\gamma_j u_{ij}^2$, where $\gamma_i$ are these eigenvalues. Bound each product between $v^{-4}$ and one to obtain both inequalities. This proves the comparison without treating the coordinate Hessian as an intrinsic curvature tensor.
\end{proof}

\begin{lemma}[Periodicity controls slope and oscillation by the Hessian]\label{lem:periodic-hessian-oscillation}
Let $d$ be the diameter of the flat torus. For $u\in C^2(\mathbb T^n)$,
\[
 \|Du\|_\infty\leq d\|D^2u\|_\infty,\qquad
 \max u-\min u\leq d\|Du\|_\infty
                  \leq d^2\|D^2u\|_\infty.
\]
Here the global parallel frame identifies tangent vectors at different points.
\end{lemma}
\begin{proof}
Every component of $Du$ has zero integral on the torus, by periodic integration by parts against a constant parallel vector. For points $x,y$, integrate the derivative of this vector-valued gradient along a shortest torus geodesic to get $|Du(x)-Du(y)|\leq d\|D^2u\|_\infty$. Average in $y$ and use the zero vector average to obtain the first estimate. Integrate $Du$ along a shortest geodesic between a maximum and a minimum for the second. A constant nonzero gradient on the Euclidean lift would not be periodic, which is precisely the obstruction removed by the averaging step.
\end{proof}

\begin{proposition}[Height intervals are nested]\label{prop:periodic-height-intervals}
For a smooth periodic graphical flow, $\max_xu(x,t)$ is nonincreasing and $\min_xu(x,t)$ is nondecreasing. Hence its height remains in the initial interval at every smooth time.
\end{proposition}
\begin{proof}
At a spatial maximum, the graph equation has $Du=0$ and gives $u_t=\Delta u\leq0$. For a fully time-dependent argument, apply the first-interior-maximum test to $u-\max u(\cdot,s)-\eta(t-s)$ on a compact time slab, using the positive definite graph coefficient matrix, and send $\eta\downarrow0$. Apply the same reasoning to $-u$. Starting at any time $s$ proves nesting, not only a bound by the initial interval. This is the compact analogue of the height estimates in \cite{EH89}, Section 2.
\end{proof}

\begin{theorem}[Conditional convergence to one horizontal slice]\label{thm:periodic-graph-conditional-convergence}
Let a periodic graphical mean curvature flow be smooth for all $t\geq0$, and put $V_0=\max\sqrt{1+|Du(\cdot,0)|^2}$. There is a constant $c_\infty$ between its initial minimum and maximum such that
\[
 \|u(\cdot,t)-c_\infty\|_\infty+
 \|Du(\cdot,t)\|_\infty+\|D^2u(\cdot,t)\|_\infty
 \leq C t^{-1/2}\quad(t>0).
\]
The constant depends only on $V_0$ and the flat torus. This is a consequence of the slope-based curvature estimate of \cite{EH89}, Proposition 4.4, with periodicity used to control the height oscillation.
\end{theorem}
\begin{proof}
Proposition \ref{prop:periodic-graph-slope-bound} bounds $v$ by $V_0$. Theorem \ref{thm:periodic-graph-curvature-smoothing} and Lemma \ref{lem:graph-curvature-hessian-comparison} give
\[
 \|D^2u(\cdot,t)\|_\infty
 \leq V_0^3\sqrt{\frac{V_0^2-1}{2t}}.
\]
Lemma \ref{lem:periodic-hessian-oscillation} gives the same rate for the gradient and oscillation. The nested closed height intervals in Proposition \ref{prop:periodic-height-intervals} are nonempty, bounded, and have length tending to zero. Their intersection is therefore one number $c_\infty$. At each time the distance of any height to this number is at most that interval's length. Combining these estimates proves the assertion. If $V_0=1$, the solution is already a constant and the estimates have zero right-hand side. The proof needs no conservation of the ordinary spatial average.
\end{proof}

\begin{proposition}[Graphical flow is an area gradient flow with variable mobility]\label{prop:graph-area-mobility}
For a smooth periodic function define its graph area
$\mathcal A(u)=\int_{\mathbb T^n}\sqrt{1+|Du|^2}\,dx$.
Its first variation and the graphical flow equation are
\[
 D\mathcal A(u)[\phi]=\int H\phi\,dx,
 \qquad u_t=-vH.
\]
Thus the flow is the negative gradient for the positive inner product
$\langle\phi,\eta\rangle_u=\int\phi\eta/v\,dx$.
Along it,
\[
 \frac d{dt}\mathcal A(u(t))=-\int H^2v\,dx
                              =-\int u_t^2/v\,dx.
\]
\end{proposition}
\begin{proof}
Differentiate under the area integral and integrate by parts on the torus:
$D\mathcal A(u)[\phi]=\int Du\cdot D\phi/v=-\int\operatorname{div}(Du/v)\phi=\int H\phi$.
The gradient with respect to the stated inner product is $vH$. Its negative is exactly the graphical equation. Substitution of $u_t=-vH$ in the first variation proves the dissipation identity; it agrees with Proposition \ref{prop:mcf-area-dissipation} because $d\mu=v\,dx$. The variable factor $v$ explains why an ordinary unweighted spatial mean is not automatically conserved.
\end{proof}

\begin{lemma}[An area-excess dissipation inequality]\label{lem:periodic-area-dissipation-gap}
Let $u$ be any smooth function on the flat torus with $1\leq v\leq V$, and set
\[
 \mathcal F(u)=\mathcal A(u)-\operatorname{Vol}(\mathbb T^n),\qquad
 \lambda_1=\inf_{\substack{f\in H^1,\ \int f=0\\f\ne0}}
            \frac{\int|Df|^2}{\int f^2}>0.
\]
Then
\[
 \frac1{V+1}\int|Du|^2\leq\mathcal F(u)\leq\frac12\int|Du|^2,
 \qquad
 \int H^2v\,dx\geq\frac{2\lambda_1}{V^2}\mathcal F(u).
\]
\end{lemma}
\begin{proof}
Use $v-1=|Du|^2/(v+1)$ to obtain the two-sided area comparison. If $Du=0$, the final estimate is immediate. Otherwise integration by parts gives
\[
 \int H(u-\bar u)\,dx=\int\frac{|Du|^2}{v}\,dx
                         \geq V^{-1}\int|Du|^2.
\]
Cauchy--Schwarz and Corollary \ref{cor:closed-manifold-poincare} bound its square above by
$\lambda_1^{-1}\int H^2\int|Du|^2$.
Cancel the positive gradient integral to get
$\int H^2\geq\lambda_1V^{-2}\int|Du|^2$.
Since $v\geq1$ and $\int|Du|^2\geq2\mathcal F$, this proves the dissipation bound. No critical-point assumption was made for this inequality.
\end{proof}

\begin{theorem}[Exponential decay of periodic graph area and gradient energy]\label{thm:periodic-area-exponential-decay}
Under Theorem \ref{thm:periodic-graph-conditional-convergence}, put $\gamma=2\lambda_1/V_0^2$. Then
\[
 \mathcal F(u(t))\leq\mathcal F(u(0))e^{-\gamma t},\qquad
 \int|Du(\cdot,t)|^2\leq(V_0+1)\mathcal F(u(0))e^{-\gamma t}.
\]
The curvature estimate of \cite{EH89} supplies the slope-controlled setting; the quantitative energy decay follows here from the graph variation and the torus Poincare inequality.
\end{theorem}
\begin{proof}
The slope bound permits $V=V_0$ in Lemma \ref{lem:periodic-area-dissipation-gap} at every time. Proposition \ref{prop:graph-area-mobility} gives
$\mathcal F'(t)\leq-\gamma\mathcal F(t)$.
Multiplication by $e^{\gamma t}$ and integration prove the first inequality. The lower area comparison from the lemma yields the second. Both inequalities hold on every existing smooth time interval; the all-time hypothesis is needed only when taking a large-time limit.
\end{proof}

\begin{lemma}[A Lipschitz-to-uniform interpolation estimate on the torus]\label{lem:torus-lipschitz-l2-interpolation}
There is a constant $C$ depending only on the flat torus such that every vector-valued Lipschitz function $f$ with Lipschitz constant at most $L\geq0$ satisfies
\[
 \|f\|_\infty\leq C\left(\|f\|_2+
                      L^{n/(n+2)}\|f\|_2^{2/(n+2)}\right).
\]
The Euclidean norm is used in the target; $C$ does not depend on its dimension.
\end{lemma}
\begin{proof}
Choose a fixed $r_0>0$ smaller than the injectivity radius, so every ball of radius $r\leq r_0$ has volume $\omega_nr^n$. Let $M=\max|f|$, attained at $x_0$. If $M=0$ there is nothing to prove. If $L=0$, the function is constant and the first term suffices. Otherwise set $r=\min\{r_0,M/(2L)\}$. The triangle inequality and the Lipschitz bound give $|f|\geq M/2$ on $B_r(x_0)$. Hence
$\|f\|_2^2\geq(\omega_n/4)M^2r^n$.
If $r=r_0$, this bounds $M$ by a fixed multiple of $\|f\|_2$. In the other case it gives $M^{n+2}\leq C L^n\|f\|_2^2$. Taking the appropriate root in each case proves the claim. Only the norm triangle inequality was used in the target, so the dimension of that target does not enter the constant.
\end{proof}

\begin{corollary}[Exponential convergence of height and slope]\label{cor:periodic-graph-exponential-c1}
In the setting of Theorem \ref{thm:periodic-graph-conditional-convergence}, with the same $c_\infty$ and $\gamma$, there is $C$ depending on the initial slope, area excess and torus such that
\[
 \|u(\cdot,t)-c_\infty\|_\infty+\|Du(\cdot,t)\|_\infty
 \leq C e^{-\gamma t/(n+2)}\qquad(t\geq1).
\]
The argument here gives exponential convergence in $C^1$ and the preceding polynomial bound in $C^2$; it does not assert uniform estimates of all derivatives without further parabolic input.
\end{corollary}
\begin{proof}
The Hessian estimate in the proof of Theorem \ref{thm:periodic-graph-conditional-convergence} gives one uniform bound $L$ for the Lipschitz constant of $Du$ for $t\geq1$. Theorem \ref{thm:periodic-area-exponential-decay} gives $\|Du\|_2\leq C e^{-\gamma t/2}$. Apply Lemma \ref{lem:torus-lipschitz-l2-interpolation} to this vector-valued gradient. Its two terms decay at rates $\gamma/2$ and $\gamma/(n+2)$, respectively, so the slower rate bounds their sum. The height interval contains $c_\infty$, and its length is at most $d\|Du\|_\infty$ by Lemma \ref{lem:periodic-hessian-oscillation}. This gives the height estimate. When the initial area excess is zero, the solution is a constant and the conclusion is immediate.
\end{proof}

\begin{corollary}[The total curvature dissipation equals the initial area excess]\label{cor:periodic-total-area-dissipation}
For an all-time smooth periodic graphical flow,
\[
 \int_0^\infty\int_{M_t}H^2\,d\mu_t\,dt
       =\mathcal A(u(0))-\operatorname{Vol}(\mathbb T^n).
\]
More generally, the remaining dissipation after a time $s$ is exactly $\mathcal F(u(s))$ and is bounded by $\mathcal F(u(0))e^{-\gamma s}$.
\end{corollary}
\begin{proof}
Integrate Proposition \ref{prop:graph-area-mobility} from $s$ to a finite $T$. The right-hand side is $\mathcal A(u(s))-\mathcal A(u(T))$. Theorem \ref{thm:periodic-area-exponential-decay} implies that the latter area tends to the base torus volume as $T\to\infty$. Monotone convergence for the nonnegative time integrand gives the exact identity. The exponential bound is the same theorem evaluated at $s$. This is a finite total dissipation statement, not a pointwise estimate for $H$ at every late time.
\end{proof}
